% 职场.tex
% Auto-generated from 感悟论语.odt by tools/odt_to_tex.py

\setpartplant{artwork/plum.png}
  {梅}
  {磨而不磷，涅而不缁。}
  {《论语·阳货》}
  {职场}
\part{职场}

\chapter{不存在“35岁陷阱”}

\begin{classic}
《论语·为政篇第二 2.12》子曰：「君子不器。」
\end{classic}

网上经常看到一些互联网企业清退那些大于等于35岁的员工，于是也便有了所谓“35岁陷阱”的说法。然而，一个人从知识体系的逐渐完善与纯熟，到设计思想的持续精进，再到技艺经验的不断锤炼与提升都是一辈子的事情，而且是没有上限的。企业硬生生地以35岁作为划分雇佣与清退的界限，只能说明不论企业表面上声称的如何，在事实上他们一直无情地将员工视为器物和工具；而且因企业本身的功利、短视、平庸，对这些“器物和工具”的利用也只是局限在较低的层次上，因而有意识地完全忽略和埋没了员工作为一个真正的人本来所具有的无限潜力。

对此，打工者不能因为企业的无理清退或是耳闻目睹到了相关的消息便给自己灌输那种“人到了一定年龄，学习、研究、编程就会干不动，记忆力与工作效率也比不过年轻人”的错误的潜意识。正确的态度是始终做自己的主人，持续提升、超越自我、止于至善。

在具体的方法上，则需要观察与评估我们是否已做到了问心无愧的努力。同时，也要反思我们采取的学习方法是否正确得当，否则便会陷入“学而不思则罔，思而不学则殆”的困境。再者，审视一下我们是否具备了良好的态度与性格。这包括无比的耐心与细心，无所畏惧、迎难而上的勇气，不怕麻烦地“敏于事、慎于言”。对此，《三体》一书也提到，真正的基础研究反倒需要不那么聪明（滑头）的人来搞——因为只有这样的人才能做到在科学面前踏实、严谨、谦虚、虔诚，从而保证了逻辑的缜密，也没有任何先入为主的自以为是。

拥有了以上正确的思想与态度，我们还需要注意身心、生理的健康。在学习与工作的过程中，我们需要保持高效但却放松的身体与精神状态，从而在最大程度上延长我们的学术与技术生命。例如，可以适当地将我们书写、打字的速度放慢下来，从而避免因过度用力而产生的肌肉紧张与劳损，也能使我们更好地感受、体会和记住在专注状态中的细微体验。这样会有助于培养我们的肌肉记忆，令技艺更为精湛、纯熟。

总之，企业没有拿员工当作真正的人来对待，员工也便不值得为企业卖命付出。打工者不要被企业设定的种种规则、制度、指标蒙蔽双眼，甚至是陷入怀疑、焦虑、恐慌。相反地，打工者要启发自己的心智，为自身的提升与精进不断努力，最终做自己生命与事业的主人。

\entrydate{2021.01.07}

\chapter{君子无终食之间违仁}

时至今时，仍然时常想起高中化学老师对我的教诲：“学习知识就要像牛儿看到满山的青草一样如饥似渴”。由此，亦回忆起我在高一时于一个周六严热的下午上化学奥赛课，一下子塞进了那么多的知识（第IV与V族元素）。可见，在将自己完全清空、严格按照训练规定去做的时候，能够产生多么巨大的动力。而且，总是可以把内在、外在的各种压力转变为动力。现在距离高中时代已近二十年了，我仍应保持那份初心与动力。原先是个人、家长、学校、老师、社会全方面对我产生的约束作用。现在能够起到监督作用的只能是靠自己。然而，君子贵在慎独、慎微，应从点滴和当下做起。

走上工作岗位后的研究与中学的学习不同的是：以前学生时代所学的都是已有的、纯粹经典的内容。那时的我不光记忆力好，还经过了无数次的反复与巩固，因而可以做到熟稔于心，使用起来得心应手、行云流水。而我现在的状况则是基础知识不甚牢固，新知识太多，记忆力亦变差、前学后忘。不仅如此，由于工作（为他人服务）的需要，无法长期专注于一个领域展开持久深入的研究。对于我当下所做的工作，由于已经进入了前沿性的研究层次，每走一步都看似希望渺茫，步履维艰。但正如论语所说：“君子无终食之间违仁，造次必于是，颠沛必于是。”无论难易，都要保持那种如饥似渴的学习与研究状态。

\begin{classic}
《论语·里仁篇第四 4.5》 子曰：「富与贵，是人之所欲也，不以其道得之，不处也。贫与贱，是人之恶也，不以其道得之，不去也。君子去仁，恶乎成名。君子无终食之间违仁，造次必于是，颠沛必于是。」
\end{classic}

\entrydate{2016.11.28}

\chapter{君子有所不为}

\begin{classic}
《论语·子张篇第十九 19.4》子夏曰：“虽小道，必有可观者焉；致远恐泥，是以君子不为也。”
\end{classic}

说一个人爱好广泛，通常是褒义的。但对于中年人来说，这其中也藏有陷阱。比如，这一句论语就警醒我们不宜在业余爱好上投入过多的精力与时间，以免影响到对人生大事业与终极目标孜孜不倦的追求。

毕竟，人的时间与精力总是有限的，爱好多且还想在其中某个方面获得出色表现，则势必要投入大量资源用于学习和反复训练。对那些想在忙碌的工作之余有更多知识与技术积累的“社畜”来说，本已将其白天8小时内最有效、最宝贵的精力都“奉献”给了雇主，下班后的专注度与耐心程度已是强弩之末。此时若再有广泛的爱好来抢占资源，情况则是雪上加霜：这会无情地减少本该在专业方面永无止境的打磨与修炼，给个人 \textbf{核心竞争力} 与 \textbf{不可替代性} 的塑造带来巨大阻碍。而没有了 \textbf{核心竞争力} 与 \textbf{不可替代性} 这两点，一个人在如今阶级固化的时代就只能沦为被资本宰杀的羔羊，或是减少个人欲望，在躺平中非暴力不合作。

另一方面，对已经步入职场的人，尤其是改弦易辙的可能性越来越少的中年人，虽然并非绝对，但大多其实早已过了培养艺术、体育、文学细胞的黄金时期。虽然成年人因为具有比孩童更健全的理性与理解力而上手更快，且即便能够将全部业余时间用于这些方面的训练，但因为时间总量的相对缺乏和缺少专业、职业的教练，也基本无法与科班出身的职业选手并驾齐驱。也许，最终的确能达到B站优质UP主自娱自乐的程度，赚得粉丝的点赞捧场，但却难以成就大创作、大事业。

幸好，两千多年前的论语在当下依旧有其适用性，能够时时提醒我们：在爱好的培养方面，君子有所不为或者不必陷得太深，中年人应学会做减法与自我收缩。

\entrydate{2022.09.30}

\chapter{选择也没有那么难}

如傅佩荣教授所说，一个人只能要么是君子，要么是小人，不存在君子与小人的迭加态。同样地，对于人生中可能面临的的种种抉择，如果根据自己为了止于至善而预先定立的原则，以二分法来看待和处理问题，那么思路便会明晰，选择也就没有那么难了。例如，究竟是选择严谨认真而是马虎浮躁？是选择正直善良还是虚伪欺骗？是选择承担更大的风险不作恶还是为了急功近利而不择手段？是愿意倾囊建百年教堂，还是盖三五十年便破败的商品房来牟利？是想要通过格物致知以明理，还是一辈子浑浑噩噩？……

由此想来，人生中许多麻烦、困惑，以及在浮华下隐匿的不安、畏惧、汗颜、悔恨都是因为走了中间路线、一再妥协（乃至没有底线）而导致的。

\begin{classic}
《论语·卫灵公篇第十五 15.2》 在陈绝粮。从者病，莫能兴。子路愠见曰：「君子亦有穷乎？」子曰：「君子固穷，小人穷斯滥矣。」

《论语·宪问篇第十四 14.23》 子曰：「君子上达，小人下达。」
\end{classic}

\entrydate{2017.02.17}

\chapter{见微知著}

如曾子所言，“战战兢兢，如临深渊，如履薄冰”。人生就是这么一个“提心吊胆”、历经艰险的过程。对于日常生活中的一言一行，都应持有无比谨慎、专注和认真的态度，需要无时无刻不控制、约束、调整自己随时可能放纵、失控的性情与行为。大到人生方向的选择，待人接物，不辱使命地完成职责，小到每日的洒扫起居、书写的一笔一划，都要予以留意，专注去做，用心体会与反省。真正的执行力也正是来源于这样严格的自我监督。

当我们谈及职业的长远发展、人生的宏大志向，往往容易感到茫然自失，难以把控。而这些构成人生历程，塑造性情和灵魂要素的点滴言行恰恰是伸手可及的。善始善终地将它们做好、做到位，便是一种修行。傅佩荣教授在《向庄子问道》系列讲座中提到：一个人若做到了完全的专注，在他的身上便会体现出神性。因此，即便一个人地位卑微、生活平淡，但岂可浑浑噩噩，应付差事般地度过？气质与品格，当于细微处塑造和体现。

\begin{classic}
《论语·泰伯篇第八 8.3》 曾子有疾，召门弟子曰：「启予足！启予手！诗云：『战战兢兢，如临深渊，如履薄冰。』而今而后，吾知免夫！小子！」
\end{classic}

\entrydate{2016.10.15}
